\section{Significance of the Study}
\label{sec:significance_of_the_study}

Addressing complex traffic networks is of paramount significance as it is crucial for growing urbanized areas such as Iloilo City (BLGF, 2025). Traffic congestion is already established along some of the primary roads in Iloilo City (Landoy, 2012), as cited in (Sosuan, 2014). This congestion serves as more than an inconvenience; as stated by the (Metropolitan Transportation Commission, 2022), traffic congestion increases travel time and leads to more unpredictable estimates of arrival. Currently, there are few to no advanced traffic management systems integrated into the development and planning of the traffic networks in the city (Conlu, 2022). Thus, the application of the Braess Paradox to these networks provides a novel way to tackle congestion by analyzing scenarios where adding road capacity might actually cause congestion to persist or worsen (Braess, 2005). 

The significance of this study extends to improving the overall quality of life in the city and ensuring sustained economic growth (Fattah, 2022). For the planning and development offices of the provincial and city government, the simulation model serves as a low-cost option for testing modifications to the road network before proceeding with final physical implementations (ICPDO, 2021). For the general public and commuters, effective road closure strategies can lead to more predictable and shorter travel times, reducing the daily stress associated with traffic-congested commutes (Marve, 2016).

For the local economy, an optimized transportation network helps businesses lower logistics expenses, such as delivery times, and establishes a system that can withstand the growing urbanized economic demands of the city for future investment (Onyeneke, 2018). Additionally, a more efficient traffic network leads to reduced vehicle emissions, directly contributing to the environmental sustainability and efficiency of the city  (Marve, 2016).